%% kapa: exam=1 round=23 year=2012 period=2 form=A subject=accounting subject_ko=회계학 seq=35 number=35
%% source: Q-Net 감정평가사 기출문제 (cst00309) artlSeq=5052580 file=raw/1008/5052580/2012년 제23회 감정평가사 제1차시험 기출문제(2교시 A형).pdf
%% conversion: kordoc 4.24.51 -> tools/gen1.py (auto); review=proofread-sonnet; flags=none
\begin{question}{35}
(주)한국은 제품A와 제품B를 거래처에 납품하는 업체이다. 제품은 생산과 동시에 전량 납품된다.
(주)한국은 합리적인 가격설정목적으로 주요활동을 분석하여 다음과 같은 자료를 수집하였다. \begin{center}
\small
\begin{tabular}{|c|r|c|r|}
\hline
활 동 & \makecell{활동별\\제조간접원가} & 원가동인 & \makecell{총원가\\동인 수} \\
\hline
제품설계 & ￦450,000 & \makecell{설계제품\\부품수} & 100단위 \\
\hline
생산준비 & 325,000 & 준비횟수 & 650회 \\
\hline
생산운영 & 637,500 & 기계작업시간 & 12,750시간 \\
\hline
선적준비 & 80,000 & 선적횟수 & 200회 \\
\hline
배 달 & 300,000 & \makecell{배달제품\\중량} & 60,000kg \\
\hline
\end{tabular}
\end{center}
제품A의 생산량은 1,000단위이고, 제품A의 원가동인이 설계제품 부품수 70단위, 준비횟수 286회, 기계작업시간 3,060시간, 선적횟수 100회, 배달제품 중량 21,000kg일 경우 활동기준원가계산방법을 적용하여 계산한 제품A의 단위당 제조간접원가는? \par
\begin{choices}
  \item ￦525
  \item ￦620
  \item ￦682
  \item ￦756
  \item ￦810
\end{choices}
\end{question}
\answer{4}
